\scanpage{44}\mspage{7}
canonique
\[\begin{aligned}
I_s\ \deleted{\unsure{D_i}}&\mathrel{\mathop{\longmapsto}\limits^{\varphi=\varphi_{s,\Delta}}}\deleted{\unsure{D_i\cap\Delta}}\ \text{ens. }\deleted{\text{des}}\ S_\Delta\text{ des sommets}\\
&\hspace{3em}\text{de }K\text{ sur }\Delta
\end{aligned}\tag*{(25)}\]
caractérisée par
\[\deleted{\unsure{\text{ou}}}\qquad D_i\cap\Delta=\{\varphi(\unsure{D|i|D_i})\}\]
Cette application est \uline{injective}, \deleted{et} i.e.
\[\varphi_{s,\Delta}:I_s\lhook\joinrel\longrightarrow S_\Delta\]
Comme sur $S_\Delta$ (ens des sommets d'un polygône)
il y a une structure polygônale, il y a sur
$I_s$ une structure polygonale induite, \deleted{\unsure{donc}}
soit $\deleted{\unsure{\varphi|Z}}Z_{s,\Delta}^{\unsure{\prime|1}}$. On aura alors

% slanted
\margin{Condition en les sommets $s$ tels que $\nu_s\geq3$.}
\[\begin{aligned} &Z'_{s,\Delta}\text{ est indépendant du choix de}\\
&\Delta\in I\setminus I_s,\quad\text{soit } Z'_s.
\end{aligned}\tag*{(A1)}\]
% slanted
\margin{à prouver que c'est vrai quand on part d'un $\mathcal X$ (et qu'on trouve bien le $Z_s$ \uline{déduit} de $Z_s$\ldots) \unsure{itou|idem|item} pour la variante avec $\vec\Delta$, \unsure{plus fine|plus loin}.}
\deleted{On} Considérons l'application
\[\vec I_s\longrightarrow I_s\qquad\text{de degré }2\tag*{(26)}\]
\deleted{On doit}\inserted{On aura} alors
\[\left\{\begin{aligned}
&\text{si }\deleted{\unsure{\rho|\varphi}}\,\vec u\in\vec{\unsure{Z|\mathcal Z|\mathfrak Z}}_s\inserted{\in{\unsure{\mathfrak S|\mathscr S|\mathcal C}}_{\vec I_s}},\quad\vec u\text{ passe au quotient dans (26),}\\
&\text{en }u\deleted{\unsure{'} }\in\mathfrak S_{I_s},\quad\text{et on a }u\in Z'_s.
\end{aligned}\right.\tag*{(27)}\]
C'est là une condition restrictive importante
sur la structure polygônale $\vec Z_s$ qu'on veut déterminer,
qui limite le choix à $2^{\nu_s-1}$ possibilités,
% archive sheet number: 43
